\documentclass{article}
\usepackage{amsmath, amssymb}
\begin{document}
\title{Digest: Physics of the Riemann Hypothesis}
\author{D\'aniel Schumayer and David A. W. Hutchinson (Digest)}
\maketitle

\section{Introduction}
This document digests the comprehensive review ``Physics of the Riemann Hypothesis.'' The paper bridges the Riemann zeta function $\zeta(s)$ with various physical phenomena, suggesting that the elusive proof of the Riemann Hypothesis might be encoded within the laws of physics.

\section{Classical and Quantum Mechanics}
The paper explores the Hilbert-P\'olya conjecture, which posits that the non-trivial zeros of $\zeta(s)$ correspond to the energy eigenvalues of a self-adjoint Hamiltonian operator. Trace formulas, such as Gutzwiller's and Selberg's, link the distribution of prime numbers to the unstable periodic orbits of classically chaotic systems.

\section{Random Matrix Theory and Nuclear Physics}
The nearest-neighbor spacing distribution of the non-trivial zeros of the Riemann zeta function exhibits striking similarity to the eigenvalue spacing of random matrices drawn from the Gaussian Unitary Ensemble (GUE), a theory originally developed to model the energy levels of heavy atomic nuclei.

\section{Statistical Mechanics: The Primon Gas}
A fascinating mathematical model discussed is the ``Riemann gas'' or ``primon gas.'' In this model, non-interacting bosons (primons) occupy energy levels proportional to the logarithm of prime numbers, $E_p = E_0 \ln p$. The resulting canonical partition function is exactly the Riemann zeta function:
\begin{equation}
Z_B(T) = \sum_{n=1}^\infty e^{-E_n / k_B T} = \sum_{n=1}^\infty n^{-s} = \zeta(s)
\end{equation}
where $s = E_0 / k_B T$. This framework offers a direct physical manifestation of the zeta function's pole at $s=1$, corresponding to the Hagedorn temperature limit.

\end{document}
